\bibitem{dyadic-bernoulli-sharifi}
Romyar T. Sharifi, \emph{Iwasawa theory}, Chapter 4,
Lemma 4.2.16 and Theorem 4.2.20: the CM unit index and the
relative class-number formula for an arbitrary abelian CM field.
\url{https://math.ucla.edu/~sharifi/notes/iwasawa-ch04.html}.
The retrieved primary-source snapshot is hashed in
\texttt{certificates/rh-dyadic-bernoulli-unit.json}.

\bibitem{dyadic-bernoulli-milne}
J. S. Milne, \emph{Class field theory}, version 4.03,
6 August 2020, Chapter V, Example 3.9 and the reciprocity
and existence theorems preceding it.
\url{https://www.jmilne.org/math/CourseNotes/CFT.pdf}.

\bibitem{dyadic-bernoulli-source}
Levent Alp\"oge, pinned source repository \texttt{jsmorph/rh},
revision \texttt{ecd09ff9b4aaf6c08f2b5ef17dc49ca9648a77de}.
In \texttt{Development.AutomorphicArithmetic.\allowbreak
AutomorphicVarieties.LevelCohomology},
\texttt{UnramifiedCocycleIsCoboundaryOfBernoulliUnit},
lines 14--26 and 37--54, specifies the exact Bernoulli
and inertia-index hypotheses; \texttt{HcVanOf}, lines 7--19,
retains local constancy at the actual cohomological consumer.
The definition of the cocycle is in
\texttt{PrimitiveFrobeniusFixed.\allowbreak IsAdmissiblePrimeListNil},
lines 98--99, in the same \texttt{AutomorphicVarieties} directory.
The full source hashes, bounded excerpts, exact polynomial
diagnostics and independent proof review are recorded in
\texttt{certificates/rh-dyadic-bernoulli-unit.json}.
